\section*{Methods}

GFP experiments use the Sarkisyan et al. local fitness landscape, with sequence records represented by one-hot and learned protein features as specified in the Supplementary Information. Materials experiments use the Matbench experimental band-gap task, with composition features and fixed train/candidate splits. The primary GFP loops run five acquisition rounds of 100 records each, for 500 selected records per seed. The primary materials loops run five rounds of 50 records each, for 250 selected records per seed. These batch sizes match the configured active-learning budgets used throughout the paired clean, random-swap and targeted-swap comparisons. Unless otherwise noted, comparisons use 10 paired seeds with shared initial histories, shared candidate pools and matched swap budgets.

Each seed defines a full closed-loop replay rather than an independent static train/test split. At round \(r\), the surrogate is trained on the current observed history, scores the remaining candidate pool, selects a batch according to the acquisition rule and then receives true feedback for the executed records. Clean, random-swap and targeted-swap conditions share the same initial history and candidate pool within a seed, so paired contrasts isolate the effect of the binding operation. Candidate records already selected in a previous round are removed from the available pool before the next acquisition step.

Paired misbinding is constructed by exchanging measured values between records while leaving the object identifiers, features and metadata fields fixed. The threat model is a structured record-operation error--for example a plate, batch, rank or sort-key offset--rather than an adversary modifying measurement values. Target slices are selected before loop execution by support and prevalence filters, low-outcome target constraints and high-outcome donor constraints; the primary targets are GFP \texttt{pos=27} and the materials Co axis. Targeted swaps attach high donor values to low-performing target records or to records satisfying a distributed conditional trigger. The label multiset is therefore preserved exactly, but the recorded conditional relation changes. Same-budget random swaps use the same number of exchanged labels but sample pairs independently of the target or trigger, separating generic label disorder from coherent relinking. New records acquired during the loop receive their true labels, so feedback is oracle feedback from the correctly bound archive rather than continued corruption. Prespecified targeted-random contrasts are evaluated by paired sign tests when all seed-wise differences have the same sign.

Coherence sweeps keep the swap count fixed and vary the fraction of swaps aligned with the target conditional direction. The remaining swaps are sampled without the target-aligned donor-target structure. This design separates error structure from error quantity: the marginal label multiset is unchanged at every coherence level, and the number of exchanged labels is constant. Generic sorted-join and block-shift controls use the same candidate universe but differ in whether the induced relinking creates a learnable target-axis contrast.

The primary surrogate is a multilayer perceptron trained independently at each round on standardized features and standardized outcomes. It uses two ReLU hidden layers with dropout 0.05, hidden width 256 for GFP and 128 for materials, and a linear output head. Models are trained for 80 epochs with Adam, learning rate 0.001, weight decay 0.0001, training batch size 256 and gradient clipping at norm 1.0. Sensitivity analyses include TabM-mini, a transformer-style tabular model, robust-loss variants, 20\% epsilon-greedy exploration and Monte Carlo-dropout upper-confidence-bound acquisition. The primary acquisition rule is top-ranked batch selection from the current candidate pool. These alternatives test whether the binding-to-budget effect depends on a single architecture or on pure exploitation.

Primary outcome summaries report final false-axis acquisition counts, true-response shortfall relative to matched controls, concentration of proposed actions and recovery rank of the implicated axis. For binary paired contrasts, the planned nonparametric test is the paired sign test across seeds; for descriptive sweeps, results are reported as seed means and threshold traces rather than as post hoc model fits. Control thresholds for quarantine are calibrated from clean and same-budget random-swap traces and then applied to targeted traces without using targeted outcomes to set the threshold.

Online quarantine is evaluated by rerunning the loop at each round. The surrogate first proposes a ranked batch; the system then computes candidate-normalized concentration for the monitored slice, quarantines over-concentrated records before execution, refills from the next ranked non-quarantined candidates and trains the next round only on executed feedback. Feedback-conflict triage is computed after true feedback returns by ranking axes according to selected-budget concentration weighted by true-response deficit. Primary endpoints are false-axis acquisition, same-budget true-response shortfall, trace-concentration stop-loss and feedback-conflict axis recovery.

External BEAR, CAMEO and SAMPLE analyses are retrospective stress replays on archived public measurements. They test whether the same mechanism and stop-loss signals remain interpretable on real measurement records, and are not audits of corruption in the original campaigns. Full model hyperparameters, seeds, thresholds, replay construction and figure-generation details are given in the Supplementary Information and reproducibility package.
